\documentclass[11pt]{article}
\usepackage[margin=1in]{geometry}
\usepackage{amsmath,amssymb,amsthm}
\usepackage{hyperref}
\usepackage{xurl}
\let\oldus\_
\renewcommand{\_}{\oldus\allowbreak}

\newtheorem{theorem}{Theorem}
\newtheorem{lemma}[theorem]{Lemma}
\theoremstyle{definition}
\newtheorem{definition}[theorem]{Definition}
\theoremstyle{remark}
\newtheorem{remark}[theorem]{Remark}

\newcommand{\C}{\mathbb{C}}
\title{A disproof of Conjecture 00000000982\\[2pt]
\large Fixed points of the Berezin transform on Fock space are neither radial nor at most two}
\date{}

\begin{document}
\sloppy
\maketitle

\section{The conjecture and its reading}

Conjecture 00000000982 reads: \emph{Definition: The Berezin transform. Conjecture: Fixed points of the
Berezin transform on Fock space are radial functions, and the cardinality of the fixed-point set is at
most 2 (fixed-point rigidity).}

It is a conjunction of two clauses: (A) every fixed point is radial, i.e.\ depends only on $|z|$;
(B) the fixed-point set has cardinality at most $2$. We use the classical Fock space of one complex
variable. The phrase ``on Fock space'' leaves the class of admissible symbols open, so we treat two classes:
\begin{itemize}
\item[(R1)] symbols in the Fock space $F^2$ itself;
\item[(R2)] bounded symbols, the standard domain of the Berezin transform of a function.
\end{itemize}
We show that (A) and (B) both fail under (R1), and (B) fails under (R2). Under (R2), clause (A) is true,
because bounded fixed points are constant (see Section~\ref{sec:lit}). The conjunction therefore still
fails under (R2).

\section{Definitions}

\begin{definition}[Fock space]
Let $dA$ be Lebesgue area measure on $\C$ and $d\lambda(w)=\pi^{-1}e^{-|w|^2}\,dA(w)$. The Fock space
$F^2$ is the space of entire functions $f$ with $\int_\C |f|^2\,d\lambda<\infty$. Its reproducing kernel
is $K(w,z)=e^{w\bar z}$, with $K(z,z)=e^{|z|^2}$. The normalized kernel is
$k_z(w)=K(w,z)/\sqrt{K(z,z)}=e^{w\bar z-|z|^2/2}$.
\end{definition}

This agrees with the Segal--Bargmann space as described on Wikipedia, where $F^2$ is ``the space of
holomorphic functions F in n complex variables satisfying the square-integrability condition''
$\pi^{-n}\int|F(z)|^2\exp(-|z|^2)\,dz<\infty$. There the kernel functions are $F_a(z)=\exp(\bar a\cdot z)$,
with $\|F_a\|^2=\exp(|a|^2)$, so $k_z=F_z/\|F_z\|$.

\begin{definition}[Berezin transform]
For a symbol $f:\C\to\C$, $(Bf)(z)=\int_\C f(w)\,|k_z(w)|^2\,d\lambda(w)$. For bounded $f$ this equals
$\langle T_f k_z,k_z\rangle$, where $T_f$ is the Toeplitz operator. This is the same pattern as the
Bergman-disk formula $(Bf)(z)=\int_D \frac{(1-|z|^2)^2}{|1-z\bar w|^4}f(w)\,dA(w)$ quoted on Wikipedia.
\end{definition}

A function $f$ is \emph{radial} if $|z|=|w|$ implies $f(z)=f(w)$. In Lean, the integral is Mathlib's
Bochner integral with respect to the Lebesgue measure \texttt{volume} on $\C$. A non-integrable integrand
has integral $0$ there. This convention does not create spurious fixed points in the classes we use.
Integrability of $w\mapsto f(w)|k_z(w)|^2$ depends on $z$, and for the classes in the statement it
holds at every $z$. For $f$ in the Fock space, Cauchy--Schwarz bounds $\int|f||k_z|^2\,d\lambda$ by
$\|f\|\,\| |k_z|^2\|$, which is finite. For a bounded symbol, $|f||k_z|^2\le\|f\|_\infty|k_z|^2$ and
$|k_z|^2\,d\lambda$ is a probability measure. A bounded fixed point $f=Bf$ is automatically
continuous, hence measurable, because $Bf$ is a Gaussian convolution. In any case, the disproof only
uses the constants and the monomials $z^k$, and for these every integral is shown to be integrable.

\section{The proof}

\begin{lemma}[Gaussian convolution form]\label{lem:conv}
$|k_z(w)|^2e^{-|w|^2}=e^{-|w-z|^2}$. Consequently $(Bf)(z)=\int_\C f(z+u)\,G(u)\,dA(u)$, where
$G(u)=\pi^{-1}e^{-|u|^2}$.
\end{lemma}
\begin{proof}
We have $|k_z(w)|^2=e^{2\,\mathrm{Re}(w\bar z)-|z|^2}$ and $|w-z|^2=|w|^2+|z|^2-2\,\mathrm{Re}(w\bar z)$.
Then substitute $w=z+u$; Lebesgue measure is translation invariant.
\end{proof}

\begin{lemma}[Gaussian moments]\label{lem:mom}
(i) $\int_\C G\,dA=1$.\\
(ii) $u\mapsto |u|^k e^{-|u|^2}$ is integrable for every $k\ge0$.\\
(iii) $\int_\C u^jG(u)\,dA(u)=0$ for every $j\ge1$.
\end{lemma}
\begin{proof}
(i) Mathlib's $\int_V e^{-b\|v\|^2}=(\pi/b)^{\dim V/2}$ with $V=\C$ gives $\int_\C e^{-|u|^2}\,dA=\pi$.

(ii) We have $x^k\le k!\,e^{x}$ and $x-x^2\le \tfrac12-\tfrac12x^2$. Hence
$|u|^ke^{-|u|^2}\le k!\,e^{1/2}e^{-|u|^2/2}$, and the right side is integrable.

(iii) Let $\omega=e^{i\pi/j}$. Rotation $u\mapsto\omega u$ is a linear isometry of $\C=\mathbb R^2$, so it
preserves Lebesgue measure, and $G(\omega u)=G(u)$. Hence $I=\int u^jG=\int(\omega u)^jG(\omega u)=\omega^jI=-I$,
so $I=0$.
\end{proof}

\begin{theorem}\label{thm:fixed}
(a) $B$ fixes every constant: $B(c)=c$.\\
(b) $B$ fixes every monomial: $B(w^k)=w^k$ for all $k\ge0$.\\
(c) Every constant and every monomial $w^k$ lies in $F^2$.\\
(d) $w\mapsto w$ is not radial.\\
(e) The monomials are linearly independent.
\end{theorem}
\begin{proof}
(a) Use Lemma~\ref{lem:conv} and Lemma~\ref{lem:mom}(i).

(b) By the binomial theorem, $(z+u)^k=\sum_m\binom km z^m u^{k-m}$. Each term is integrable against $G$
by Lemma~\ref{lem:mom}(ii). By Lemma~\ref{lem:mom}(iii), every term with $k-m\ge1$ integrates to $0$.
The term $m=k$ integrates to $z^k$.

(c) Constants and monomials are entire. Moreover $|w^k|^2\pi^{-1}e^{-|w|^2}=\pi^{-1}|w|^{2k}e^{-|w|^2}$ is
integrable by Lemma~\ref{lem:mom}(ii).

(d) $|1|=|i|$, but $1\ne i$.

(e) The map $p\mapsto(z\mapsto p(z))$ is injective on $\C[X]$, because a polynomial vanishing on $\C$ is
$0$. It sends the monomial basis to the functions $w\mapsto w^k$.
\end{proof}

\begin{theorem}[Main theorem]
The conjecture is false.
\begin{itemize}
\item Under (R1), $w\mapsto w$ is a non-radial fixed point in $F^2$, so (A) fails. The fixed-point set
contains every constant, so it is infinite and (B) fails. It also contains infinitely many linearly
independent functions $w^k$, so (B) fails even if ``cardinality'' is read as ``dimension''.
\item Under (R2), the fixed-point set contains every constant, so it is infinite and (B) fails.
\end{itemize}
\end{theorem}

\section{Literature check}\label{sec:lit}
The result is consistent with the literature. Casseli (arXiv:1903.05364) states: ``We show that an
$L^p$ function, $p\in[1,+\infty]$, with respect to the Lebesgue measure is
invariant under this transformation if and only if it is harmonic. From this we deduce that the only
bounded fixed points of the Berezin transform of polyanalytic Fock spaces are constant functions.'' The
same paper notes: ``For the classical Fock space , a Lipschitz estimate for the Berezin transform of a
bounded function, together with a semigroup property, shows that a bounded fixed point of the Berezin
transform must be constant''. Asghari, Cuckovic and Sahutoglu (arXiv:2508.13115) state: ``It is known
that the Berezin transform is well-defined on the polynomials in $z$ and $\overline{z}$.'' On the polynomial domain, $w\mapsto w$ is again a non-radial fixed point. The proof
above does not depend on these papers.

\section{Relation to the earlier submission}
An earlier submission by earthking11 was merged and then removed in the re-audit (commit 541cf4fb). The
re-audit's reason was: ``no Berezin/Fock objects; load-bearing theorem prose-only''. Its Lean used only
\texttt{import Std}. It modelled $z$ on the integer lattice and proved only that $(1,0)\ne(0,1)$ and that
$1,z,z^2$ are distinct. The identity $B=e^{\Delta/4}$ and the fixed-point property appeared only in prose.
The present submission defines in Lean the Fock space, its normalized kernel and the Berezin transform as
an actual Lebesgue integral over $\C$. It proves the Gaussian convolution form, $B(c)=c$ and
$B(w^k)=w^k$ from Mathlib's Gaussian integral, as well as $F^2$ membership, non-radiality and the
infinitude of the fixed-point set.

\section{Lean formalization}
The file is \texttt{lean/Conjecture982/Basic.lean} (namespace \texttt{C982}), and it imports only Mathlib.

\begin{itemize}
\item Definitions: \texttt{gaussDensity}, \texttt{InFock}, \texttt{fockKernel}, \texttt{normKernel},
\texttt{berezin}, \texttt{IsRadial}, \texttt{fockFixed}, \texttt{boundedFixed}.
\item Lemmas:
  \begin{itemize}
  \item \texttt{normKernel\_eq} states $k_z=K(\cdot,z)e^{-|z|^2/2}$;
  \item \texttt{kernel\_weight} and \texttt{berezin\_eq\_conv} give Lemma~\ref{lem:conv};
  \item \texttt{integral\_G}, \texttt{integrable\_moment} and \texttt{moment\_zero} give
    Lemma~\ref{lem:mom};
  \item \texttt{berezin\_const}, \texttt{berezin\_pow}, \texttt{const\_inFock}, \texttt{pow\_inFock},
    \texttt{id\_not\_radial}, \texttt{monomials\_linearIndependent} and \texttt{fixed\_infinite} give
    Theorem~\ref{thm:fixed} and the infinitude of the fixed-point set.
  \end{itemize}
\item Main theorem \texttt{conjecture\_982\_false}. It states that $(w\mapsto w)\in$ \texttt{fockFixed} and
is not radial, and that every $w^k\in$ \texttt{fockFixed}. It also states that the monomials are linearly
independent, that \texttt{fockFixed} and \texttt{boundedFixed} are infinite, and that
$\neg((\forall f\in\texttt{fockFixed},\ \mathrm{radial}\ f)\wedge |\texttt{fockFixed}|\le2)$ and
$\neg(|\texttt{boundedFixed}|\le 2)$, with cardinalities measured by \texttt{Set.encard}.
\item Axiom audit: \texttt{\#print axioms C982.conjecture\_982\_false} reports \texttt{[propext,
Classical.choice, Quot.sound]}. There is no \texttt{sorry} and no \texttt{native\_decide}.
\end{itemize}

\end{document}
